% Part: incompleteness
% Chapter: arithmetization-syntax
% Section: proofs-in-nd

\documentclass[../../../include/open-logic-section]{subfiles}

\begin{document}

\olfileid{inc}{art}{pnd}
\olsection{\usetoken{P}{derivation} في الاستنباط الطبيعي}

\begin{explain}
حوسبة ال!!{derivation}s تبدأ بتمثيل
ال!!{derivation}s بالأعداد. ولما كانت
ال!!{derivation}s أشجارًا من !!{formula}s،
ولكل استدلال وسم أو وسمان، ناسب أن
يكون التمثيل عوديًا. فنمثل
ال!!{derivation} بصف يتضمن عدد
ال!!{derivation}s الجزئية المباشرة
التي تنتهي بمقدمات آخر استدلال،
ورموز هذه ال!!{derivation}s الجزئية،
وال!!{formula} الختامية، ووسم التفريغ،
ثم عددًا يعين نوع آخر استدلال.
\end{explain}

\begin{defn}
  لكل $\delta$ وهو !!a{derivation} في الاستنباط الطبيعي،
  نحدّ $\Gn{\delta}$ استقرائيًا بحسب الحالات الآتية:
  \begin{enumerate}
  \item 
    إن اقتصر $\delta$ على فرض~$!A$، كان $\Gn{\delta}$
    هو $\tuple{\OLMathNumeral{0},\Gn{!A},\OLId{latin-ordinary}{n}}$. ونجعل~$\OLId{latin-ordinary}{n}$ هو $\OLMathNumeral{0}$
    إن كان الفرض !!{undischarged}، ورقم وسمه فيما عدا ذلك.
  \item 
    وإن انتهى $\delta$ باستدلال بلا مقدمات، أو من مقدمة،
    أو مقدمتين، أو ثلاث، كانت قيمة $\Gn{\delta}$
    \begin{align*}
      & \tuple{\OLMathNumeral{0}, \Gn{!A}, \OLId{latin-ordinary}{n}, \OLId{latin-ordinary}{k}}, \\
      & \tuple{\OLMathNumeral{1}, \Gn{\delta_\OLMathNumeral{1}}, \Gn{!A}, \OLId{latin-ordinary}{n}, \OLId{latin-ordinary}{k}}, \\
      & \tuple{\OLMathNumeral{2}, \Gn{\delta_\OLMathNumeral{1}}, \Gn{\delta_\OLMathNumeral{2}}, \Gn{!A}, \OLId{latin-ordinary}{n}, \OLId{latin-ordinary}{k}}, \text{ أو}\\
      & \tuple{\OLMathNumeral{3}, \Gn{\delta_\OLMathNumeral{1}}, \Gn{\delta_\OLMathNumeral{2}}, \Gn{\delta_\OLMathNumeral{3}}, \Gn{!A}, \OLId{latin-ordinary}{n},
        \OLId{latin-ordinary}{k}},
    \end{align*}
    على الترتيب. وفي هذه الصفوف $\delta_\OLMathNumeral{1}$ و$\delta_\OLMathNumeral{2}$
    و$\delta_\OLMathNumeral{3}$ هي ال!!{derivation}s الجزئية التي تنتهي
    بمقدمات آخر استدلال في $\delta$؛ و$!A$ نتيجة
    آخر استدلال في~$\delta$؛ و$\OLId{latin-ordinary}{n}$ وسم تفريغه،
    ونجعله $\OLMathNumeral{0}$ إن لم يفرّغ فرضًا. وأما~$\OLId{latin-ordinary}{k}$
    فتعينه القاعدة الأخيرة وفق هذا الجدول.

    \begin{tabular}{lccccccc}
      \text{القاعدة:} & \Intro{\land} & \Elim{\land} & \Intro{\lor} & \Elim{\lor} \\
      $\OLId{latin-ordinary}{k}$: & \OLClassicalDigits{1} & \OLClassicalDigits{2} & \OLClassicalDigits{3} & \OLClassicalDigits{4}  \\[1.5ex]
      \text{القاعدة:} &  \Intro{\lif} & \Elim{\lif} & \Intro{\lnot} & \Elim{\lnot} \\
      $\OLId{latin-ordinary}{k}$: & \OLClassicalDigits{5} & \OLClassicalDigits{6} & \OLClassicalDigits{7} & \OLClassicalDigits{8} \\[1.5ex]
      \text{القاعدة:}  &  \FalseInt & \FalseCl & \Intro{\lforall} & \Elim{\lforall} \\
      $\OLId{latin-ordinary}{k}$: & \OLClassicalDigits{9} & \OLClassicalDigits{10} & \OLClassicalDigits{11} & \OLClassicalDigits{12} \\[1.5ex]
      \text{القاعدة:} & \Intro{\lexists} & \Elim{\lexists} & \Intro{\eq} & \Elim{\eq} \\
      $\OLId{latin-ordinary}{k}$: & \OLClassicalDigits{13} & \OLClassicalDigits{14} & \OLClassicalDigits{15} & \OLClassicalDigits{16} 
    \end{tabular}
  \end{enumerate}
\end{defn}

\begin{ex}
  خذ هذا ال!!{derivation} اليسير:
  \begin{prooftree}
    \AxiomC{$\Discharge{!A \land !B}{\OLMathNumeral{1}}$}
    \RightLabel{\Elim{\land}}
    \UnaryInfC{$!A$}
    \DischargeRule{\Intro{\lif}}{1}
    \UnaryInfC{$(!A \land !B) \lif !A$}
  \end{prooftree}
  فرمز الفرض هو $\OLId{latin-ordinary}{d}_\OLMathNumeral{0}=\tuple{\OLMathNumeral{0},\Gn{!A \land !B},\OLMathNumeral{1}}$.
  ورمز ال!!{derivation} المنتهي بنتيجة~$\Elim{\land}$
  هو $\OLId{latin-ordinary}{d}_\OLMathNumeral{1}=\tuple{\OLMathNumeral{1},\OLId{latin-ordinary}{d}_\OLMathNumeral{0},\Gn{!A},\OLMathNumeral{0},\OLMathNumeral{2}}$.
  وفيه $\OLMathNumeral{1}$ لأن $\Elim{\land}$ أحادية المقدمة،
  ثم عدد غودل للنتيجة~$!A$، ثم $\OLMathNumeral{0}$ لعدم
  تفريغ أي فرض، ثم $\OLMathNumeral{2}$ رمز القاعدة $\Elim{\land}$.
  فيكون رمز ال!!{derivation} كله
  $\tuple{\OLMathNumeral{1},\OLId{latin-ordinary}{d}_\OLMathNumeral{1},\Gn{((!A \land !B) \lif !A)},\OLMathNumeral{1},\OLMathNumeral{5}}$،
  أي
  \[
  \tuple{\OLMathNumeral{1}, \tuple{\OLMathNumeral{1}, \tuple{\OLMathNumeral{0}, \Gn{(!A \land !B)}, \OLMathNumeral{1}}, \Gn{!A}, \OLMathNumeral{0}, \OLMathNumeral{2}},
    \Gn{((!A \land !B) \lif !A)}, \OLMathNumeral{1}, \OLMathNumeral{5}}.
  \]
\end{ex}

\begin{explain}
وبعد تعيين تمثيل ال!!{derivation}s،
نحتاج إلى إثبات أن معالجة أعداد غودل
لهذه ال!!{derivation}s وخواصها وعلاقاتها
الأساسية عودية بدائية. فمن ذلك عمليات
يسيرة: إذا أعطي عدد غودل~$\OLId{latin-ordinary}{d}$
ل!!a{derivation}، فإن
$\fn{EndFmla}(\OLId{latin-ordinary}{d})=(\OLId{latin-ordinary}{d})_{(\OLId{latin-ordinary}{d})_\OLMathNumeral{0}+\OLMathNumeral{1}}$
تعطي عدد غودل لل!!{formula} الختامية،
و$\fn{DischargeLabel}(\OLId{latin-ordinary}{d})=(\OLId{latin-ordinary}{d})_{(\OLId{latin-ordinary}{d})_\OLMathNumeral{0}+\OLMathNumeral{2}}$
تعطي وسم التفريغ،
و$\fn{LastRule}(\OLId{latin-ordinary}{d})=(\OLId{latin-ordinary}{d})_{(\OLId{latin-ordinary}{d})_\OLMathNumeral{0}+\OLMathNumeral{3}}$
تعطي رمز نوع الاستدلال الأخير.
وثمة عمليات أشق نبيّن طريقها، ولو
إجمالًا. والمراد إثبات أن «$\delta$
!!a{derivation} لـ$!A$ من $\OLId{greek-variable}{Gamma}$»
علاقة عودية بدائية في عددي غودل
لـ$\delta$ و$!A$.
\end{explain}

\begin{prop}
  كل من العلاقتين الآتيتين عودية بدائية:
  \begin{enumerate}
  \item وجود $!A$ فرضًا في~$\delta$ موسومًا بـ~$\OLId{latin-ordinary}{n}$.
  \item كون فروض $\delta$ ذات الوسم~$\OLId{latin-ordinary}{n}$ كلها من الصورة~$!A$
    (أي إمكان !!{discharge} الفرض~$!A$ باستعمال الوسم~$\OLId{latin-ordinary}{n}$ في~$\delta$).
  \end{enumerate}
\end{prop}

\begin{proof}
  نثبت عودية العلاقتين المناظرتين بدائيًا، على أعداد
  غودل لل!!{formula}s وأعداد غودل لل!!{derivation}s.
  \begin{enumerate}
  \item أما $\fn{Assum}(\OLId{latin-ordinary}{x},\OLId{latin-ordinary}{d},\OLId{latin-ordinary}{n})$ فمعناها أن $\OLId{latin-ordinary}{x}$
    عدد غودل لفرض في !!{derivation} رمزه~$\OLId{latin-ordinary}{d}$،
    ووسم الفرض $\OLId{latin-ordinary}{n}$. فيكفي أن يكون ال!!{derivation}
    ذو الرمز $\tuple{\OLMathNumeral{0},\OLId{latin-ordinary}{x},\OLId{latin-ordinary}{n}}$ !!{derivation}ًا جزئيًا
    من~$\OLId{latin-ordinary}{d}$. وترميز ال!!{derivation}s هنا من جملة
    ترميز الأشجار في \olref[cmp][rec][tre]{sec}.
    فالدالة العودية البدائية $\fn{SubtreeSeq}(\OLId{latin-ordinary}{d})$
    تسرد رموز ال!!{derivation}s الجزئية كلها
    من~$\OLId{latin-ordinary}{d}$ في متتالية طولها لا يتجاوز $\OLId{latin-ordinary}{d}$.
    ولهذا يصح التعريف
    \[
    \fn{Assum}(\OLId{latin-ordinary}{x}, \OLId{latin-ordinary}{d}, \OLId{latin-ordinary}{n}) \defiff \bexists{\OLId{latin-ordinary}{i}<\OLId{latin-ordinary}{d}}{(\fn{SubtreeSeq}(\OLId{latin-ordinary}{d}))_\OLId{latin-ordinary}{i} =
      \tuple{\OLMathNumeral{0}, \OLId{latin-ordinary}{x}, \OLId{latin-ordinary}{n}}}.
    \]
  \item وأما $\fn{Discharge}(\OLId{latin-ordinary}{x},\OLId{latin-ordinary}{d},\OLId{latin-ordinary}{n})$، أي كون
    الفروض ذات الوسم~$\OLId{latin-ordinary}{n}$ في !!{derivation} رمزه~$\OLId{latin-ordinary}{d}$
    كلها ال!!{formula} ذات الرمز~$\OLId{latin-ordinary}{x}$، فعوديتها
    البدائية ظاهرة من تكافئها مع
    $\bforall{\OLId{latin-ordinary}{y}<\OLId{latin-ordinary}{d}}{(\fn{Assum}(\OLId{latin-ordinary}{y},\OLId{latin-ordinary}{d},\OLId{latin-ordinary}{n}) \lif \OLId{latin-ordinary}{y}=\OLId{latin-ordinary}{x})}$.
  \end{enumerate}
\end{proof}

\begin{prop}
  \ollabel{prop:followsby}
  الخاصية $\fn{Correct}(\OLId{latin-ordinary}{d})$ عودية بدائية؛ ومفادها صحة
  الاستدلال الأخير في !!{derivation}~$\delta$ الذي رمزه~$\OLId{latin-ordinary}{d}$.
\end{prop}

\begin{proof}
  لكل قاعدة استدلال~$\OLId{latin-looped}{R}$ نثبت عودية
  $\fn{FollowsBy}_\OLId{latin-looped}{R}(\OLId{latin-ordinary}{d})$ البدائية.
  ومعنى $\fn{FollowsBy}_\OLId{latin-looped}{R}(\OLId{latin-ordinary}{d})$ أن $\OLId{latin-ordinary}{d}$
  رمز !!{derivation}~$\delta$، وأن
  ال!!{formula} الختامية لـ$\delta$
  تلزم من ال!!{derivation}s الجزئية
  المباشرة لـ$\delta$ بتطبيق صحيح لـ~$\OLId{latin-looped}{R}$.

  ومن أيسر الحالات $\Intro{\land}$.
  فإن ختم $\delta$ باستدلال صحيح بحسب $\Intro{\land}$،
  كانت صورته:
  \begin{prooftree}
    \AxiomC{}
    \RightLabel{$\delta_\OLMathNumeral{1}$}
    \DeduceC{$!A$}
    
    \AxiomC{}
    \RightLabel{$\delta_\OLMathNumeral{2}$}
    \DeduceC{$!B$}

    \RightLabel{\Intro\land}
    \BinaryInfC{$!A \land !B$}
  \end{prooftree}
  فيكون عدد غودل~$\OLId{latin-ordinary}{d}$ لـ$\delta$
  $\tuple{\OLMathNumeral{2},\OLId{latin-ordinary}{d}_\OLMathNumeral{1},\OLId{latin-ordinary}{d}_\OLMathNumeral{2},\Gn{(!A \land !B)},\OLMathNumeral{0},\OLId{latin-ordinary}{k}}$، وفيه
  $\fn{EndFmla}(\OLId{latin-ordinary}{d}_\OLMathNumeral{1})=\Gn{!A}$ و$\fn{EndFmla}(\OLId{latin-ordinary}{d}_\OLMathNumeral{2})=\Gn{!B}$،
  و$\OLId{latin-ordinary}{n}=\OLMathNumeral{0}$ و$\OLId{latin-ordinary}{k}=\OLMathNumeral{1}$. فتعرّف $\fn{FollowsBy}_{\Intro\land}(\OLId{latin-ordinary}{d})$ بـ
  \begin{multline*}
    (\OLId{latin-ordinary}{d})_\OLMathNumeral{0} = \OLMathNumeral{2} \land \fn{DischargeLabel}(\OLId{latin-ordinary}{d}) = \OLMathNumeral{0} \land \fn{LastRule}(\OLId{latin-ordinary}{d}) = \OLMathNumeral{1} \land {}\\
    \fn{EndFmla}(\OLId{latin-ordinary}{d}) = {}\\
    \Gn{(} \concat \fn{EndFmla}((\OLId{latin-ordinary}{d})_\OLMathNumeral{1}) \concat \Gn{\land}
    \concat \fn{EndFmla}((\OLId{latin-ordinary}{d})_\OLMathNumeral{2}) \concat \Gn{)}.
  \end{multline*}

  ومن الأمثلة اليسيرة أيضًا $\Intro\eq$.
  فلا مقدمات لها، ولذلك $(\OLId{latin-ordinary}{d})_\OLMathNumeral{0}=\OLMathNumeral{0}$،
  كما في الفروض، ولا تفريغ فيها، أي
  $\OLId{latin-ordinary}{n}=\OLMathNumeral{0}$. غير أن $!A$ لا بد أن تكون
  $\eq[\OLId{latin-ordinary}{t}][\OLId{latin-ordinary}{t}]$ لحد مغلق~$\OLId{latin-ordinary}{t}$.
  فنحصل على التعريف العودي البدائي
  \begin{multline*}
    (\OLId{latin-ordinary}{d})_\OLMathNumeral{0} = \OLMathNumeral{0} \land \fn{DischargeLabel}(\OLId{latin-ordinary}{d}) = \OLMathNumeral{0} \land {}\\
    \bexists{\OLId{latin-ordinary}{t}<\OLId{latin-ordinary}{d}}{(\fn{ClTerm}(\OLId{latin-ordinary}{t}) \land 
      \fn{EndFmla}(\OLId{latin-ordinary}{d}) = {}\\
      \Gn{{\eq}(} \concat \OLId{latin-ordinary}{t} \concat \Gn{,} \concat \OLId{latin-ordinary}{t} \concat \Gn{)})}.
  \end{multline*}

  وأما $\fn{FollowsBy}_{\Intro{\lif}}(\OLId{latin-ordinary}{d})$
  فمثال أدق. تصدق إذا وفقط إذا كانت
  ال!!{formula} الختامية لـ$\delta$
  هي $(!A\lif !B)$، وال!!{formula}
  الختامية لـ$\delta_\OLMathNumeral{1}$ هي~$!B$،
  وكانت فروض~$\delta$ الموسومة بـ$\OLId{latin-ordinary}{n}$
  كلها من الصورة~$!A$. وهذا يعبر عنه
  عوديًا بدائيًا بالشرط
  \begin{multline*}
    (\OLId{latin-ordinary}{d})_\OLMathNumeral{0} = \OLMathNumeral{1} \land {}\\
    \bexists{\OLId{latin-ordinary}{a}<\OLId{latin-ordinary}{d}}{(\fn{Discharge}(\OLId{latin-ordinary}{a}, (\OLId{latin-ordinary}{d})_\OLMathNumeral{1}, \fn{DischargeLabel}(\OLId{latin-ordinary}{d})) \land {}}\\
      \fn{EndFmla}(\OLId{latin-ordinary}{d}) = (\Gn{(} \concat \OLId{latin-ordinary}{a} \concat \Gn{\lif}
      \concat \fn{EndFmla}((\OLId{latin-ordinary}{d})_\OLMathNumeral{1}) \concat \Gn{)}))
  \end{multline*}
  (وفيه $\OLId{latin-ordinary}{a}$ عدد غودل لـ$!A$.)

  ولننظر أيضًا في $\Intro{\lexists}$.
  فصحة آخر استدلال في $\delta$ تكافئ
  وجود !!a{formula}~$!A$ وحد مغلق~$\OLId{latin-ordinary}{t}$
  و!!a{variable}~$\OLId{latin-ordinary}{x}$، بحيث تكون
  $\Subst{!A}{\OLId{latin-ordinary}{t}}{\OLId{latin-ordinary}{x}}$ ال!!{formula}
  الختامية ل!!{derivation}~$\delta_\OLMathNumeral{1}$،
  وتكون $\lexists[\OLId{latin-ordinary}{x}][!A]$ النتيجة الأخيرة.
  فشرط $\fn{FollowsBy}_{\Intro{\lexists}}(\OLId{latin-ordinary}{d})$ هو
  \begin{multline*}
    (\OLId{latin-ordinary}{d})_\OLMathNumeral{0} = \OLMathNumeral{1} \land \fn{DischargeLabel}(\OLId{latin-ordinary}{d}) = \OLMathNumeral{0} \land {} \\
    \bexists{\OLId{latin-ordinary}{a} < \OLId{latin-ordinary}{d}}{\bexists{\OLId{latin-ordinary}{x}<\OLId{latin-ordinary}{d}}{\bexists{\OLId{latin-ordinary}{t}<\OLId{latin-ordinary}{d}}{
          (\fn{ClTerm}(\OLId{latin-ordinary}{t}) \land \fn{Var}(\OLId{latin-ordinary}{x})  \land {}}}}\\
    \fn{Subst}(\OLId{latin-ordinary}{a},\OLId{latin-ordinary}{t},\OLId{latin-ordinary}{x}) = \fn{EndFmla}((\OLId{latin-ordinary}{d})_\OLMathNumeral{1}) \land
    \fn{EndFmla}(\OLId{latin-ordinary}{d}) = (\Gn{\lexists} \concat \OLId{latin-ordinary}{x} \concat \OLId{latin-ordinary}{a})).
  \end{multline*}

  ونجمع الحالات في تعريف $\fn{Correct}(\OLId{latin-ordinary}{d})$:
  \begin{multline*}
    \fn{Sent}(\fn{EndFmla}(\OLId{latin-ordinary}{d})) \land [ {}\\
    (\fn{LastRule}(\OLId{latin-ordinary}{d}) = \OLMathNumeral{1} \land
    \fn{FollowsBy}_{\Intro\land}(\OLId{latin-ordinary}{d})) \lor \dots \lor {}\\
    (\fn{LastRule}(\OLId{latin-ordinary}{d}) = \OLMathNumeral{16} \land \fn{FollowsBy}_{\Elim\eq}(\OLId{latin-ordinary}{d})) \lor {}\\
    \bexists{\OLId{latin-ordinary}{n}<\OLId{latin-ordinary}{d}}{\bexists{\OLId{latin-ordinary}{x}<\OLId{latin-ordinary}{d}}{(\OLId{latin-ordinary}{d} = \tuple{\OLMathNumeral{0}, \OLId{latin-ordinary}{x}, \OLId{latin-ordinary}{n}})}}].
  \end{multline*}
  فيشترط السطر الأول أن تكون ال!!{formula} الختامية
  لـ$\OLId{latin-ordinary}{d}$ جملة، ويخص السطر الأخير حالة اقتصار $\OLId{latin-ordinary}{d}$ على فرض.
\end{proof}

\begin{prob}
  ضع تعريفًا للخواص الآتية على مثال
  \olref[inc][art][pnd]{prop:followsby}:
  \begin{enumerate}
  \item $\fn{FollowsBy}_{\Elim{\lif}}(\OLId{latin-ordinary}{d})$,
  \item $\fn{FollowsBy}_{\Elim{\eq}}(\OLId{latin-ordinary}{d})$,
  \item $\fn{FollowsBy}_{\Elim{\lor}}(\OLId{latin-ordinary}{d})$,
  \item $\fn{FollowsBy}_{\Intro{\lforall}}(\OLId{latin-ordinary}{d})$.
  \end{enumerate}
  وعليك في الأخيرة أن تثبت عودية اختبار شرط
  المتغير المميز بدائيًا، في آخر استدلال من
  !!{derivation} رمزه~$\OLId{latin-ordinary}{d}$. أي إن المتغير
  المميز~$\OLId{latin-ordinary}{a}$ في قاعدة $\Intro{\lforall}$
  لا يرد في ال!!{formula} الختامية لـ$\OLId{latin-ordinary}{d}$
  ولا في أي فرض مفتوح فيها. ولك أن تستعمل
  المحمول العودي البدائي $\fn{OpenAssum}$
  من \olref[inc][art][pnd]{prop:openassum}.
\end{prob}

\begin{prop}
  \ollabel{prop:deriv}
  العلاقة $\fn{Deriv}(\OLId{latin-ordinary}{d})$ عودية بدائية؛ وهي تميّز
  كون $\OLId{latin-ordinary}{d}$ عدد غودل ل!!{derivation} صحيح~$\delta$.
\end{prop}

\begin{proof}
  إن !!^a{derivation}~$\delta$ صحيح متى صحت
  تطبيقات القواعد فيه كلها؛ أي متى انتهت
  ال!!{derivation}s الجزئية جميعًا باستدلال
  صحيح. ولتكن $\fn{TreeCode}(\OLId{latin-ordinary}{d})$ الخاصية العودية البدائية الدالة
  على أن~$\OLId{latin-ordinary}{d}$ رمز شجرة منتهية ذات جذر بحسب ترميز الأشجار المعتمد.
  فـ$\fn{Deriv}(\OLId{latin-ordinary}{d})$ تكافئ
  \[
  \fn{TreeCode}(\OLId{latin-ordinary}{d}) \land
  \bforall{\OLId{latin-ordinary}{i}<\len{\fn{SubtreeSeq}(\OLId{latin-ordinary}{d})}}{\fn{Correct}((\fn{SubtreeSeq}(\OLId{latin-ordinary}{d}))_\OLId{latin-ordinary}{i})}
  \]
  % تصحيح مصدر (R287-03): أُضيف حارس رمز الشجرة لئلا تحقق الشفرات الرديئة شرط الصحة صدقًا خاليًا.
\end{proof}

\begin{prop}
  \ollabel{prop:openassum} العلاقة $\fn{OpenAssum}(\OLId{latin-ordinary}{z},\OLId{latin-ordinary}{d})$ عودية بدائية؛
  وهي تعني أن $\OLId{latin-ordinary}{z}$ عدد غودل لفرض~$!A$
  !!a{undischarged} في ال!!{derivation}~$\delta$
  ذي عدد غودل~$\OLId{latin-ordinary}{d}$.
\end{prop}

\begin{proof}
  يكون ورود الفرض !!{discharged} إذا
  حمل الوسم~$\OLId{latin-ordinary}{n}$ في !!{derivation}
  جزئي من~$\delta$، وكانت قاعدته
  الأخيرة ذات وسم تفريغ~$\OLId{latin-ordinary}{n}$.
  فلكي تكون $!A$ فرضًا !!a{undischarged}
  في~$\delta$ يكفي أن يبقى ورود لها
  واحد على الأقل غير !!{discharged}
  في~$\delta$. وينبغي التمييز بين
  الورودات؛ فقد تكون في $\delta$
  ورودات لـ$!A$ منها ما هو
  !!{discharged} ومنها ما هو !!{undischarged}.

  لنأخذ متتالية $\delta_\OLMathNumeral{0}$، \dots،
  $\delta_\OLId{latin-ordinary}{k}$، أولها $\delta_\OLMathNumeral{0}=\delta$،
  وآخرها $\delta_\OLId{latin-ordinary}{k}$ هو الفرض
  $\Discharge{!A}{\OLId{latin-ordinary}{n}}$ لبعض~$\OLId{latin-ordinary}{n}$؛
  ويكون $\delta_{\OLId{latin-ordinary}{i}+\OLMathNumeral{1}}$ !!{derivation}
  جزئيًا مباشرًا من~$\delta_\OLId{latin-ordinary}{i}$.
  فإن كان $\OLId{latin-ordinary}{n}=\OLMathNumeral{0}$ فالورود غير !!{discharged}
  بحكم التعريف؛ وإن كان $\OLId{latin-ordinary}{n}>\OLMathNumeral{0}$ لم يُفرَّغ
  ما لم تختم عقدة سابقة للفرض في هذا
  المسار باستدلال وسم تفريغه~$\OLId{latin-ordinary}{n}$.

  وعلاقة كون !!{derivation} جزئيًا مباشرًا
  عودية بدائية؛ إذ
  $\fn{Subderiv}(\OLId{latin-ordinary}{d},\OLId{latin-ordinary}{d}')$ تكافئ
  $\bexists{\OLId{latin-ordinary}{j}<(\OLId{latin-ordinary}{d}')_\OLMathNumeral{0}}{\OLId{latin-ordinary}{d}=(\OLId{latin-ordinary}{d}')_{\OLId{latin-ordinary}{j}+\OLMathNumeral{1}}}$.
  ولتكن $\fn{Path}(\OLId{latin-ordinary}{s},\OLId{latin-ordinary}{d})$ العلاقة التي
  تقول إن $\OLId{latin-ordinary}{s}$ رمز متتالية غير خالية،
  أولها~$\OLId{latin-ordinary}{d}$، وكل تالٍ فيها !!{derivation}
  جزئي مباشر من سابقه؛ أي
  \[
    \fn{SeqCode}(\OLId{latin-ordinary}{s}) \land \len{\OLId{latin-ordinary}{s}}>\OLMathNumeral{0} \land (\OLId{latin-ordinary}{s})_\OLMathNumeral{0}=\OLId{latin-ordinary}{d} \land
    \bforall{\OLId{latin-ordinary}{i}<(\len{\OLId{latin-ordinary}{s}}\tsub \OLMathNumeral{1})}{\fn{Subderiv}((\OLId{latin-ordinary}{s})_{\OLId{latin-ordinary}{i}+\OLMathNumeral{1}},(\OLId{latin-ordinary}{s})_\OLId{latin-ordinary}{i})}.
  \]
  فهي عودية بدائية. ورمز كل شجرة جزئية
  مباشرة دون رمز الشجرة الحاوية لها في
  ترميز المتتاليات المستعمل؛ فلا يجاوز
  طول مسار بادئ عند~$\OLId{latin-ordinary}{d}$ العدد $\OLId{latin-ordinary}{d}+\OLMathNumeral{1}$،
  ولا يجاوز عنصر منه~$\OLId{latin-ordinary}{d}$. ولذلك يحد
  $\fn{sequenceBound}(\OLId{latin-ordinary}{d},\OLId{latin-ordinary}{d}+\OLMathNumeral{1})$ رمز كل
  مسار كهذا. ومن ثم نعرّف
  $\fn{OpenAssum}(\OLId{latin-ordinary}{z},\OLId{latin-ordinary}{d})$ بـ
  \begin{multline*}
    \fn{TreeCode}(\OLId{latin-ordinary}{d}) \land {}\\
    \bexists{\OLId{latin-ordinary}{s}\leq\fn{sequenceBound}(\OLId{latin-ordinary}{d},\OLId{latin-ordinary}{d}+\OLMathNumeral{1})}{(\fn{Path}(\OLId{latin-ordinary}{s},\OLId{latin-ordinary}{d}) \land {}} \\
    \bexists{\OLId{latin-ordinary}{n}<\OLId{latin-ordinary}{d}}{((\OLId{latin-ordinary}{s})_{\len{\OLId{latin-ordinary}{s}} \tsub \OLMathNumeral{1}} = \tuple{\OLMathNumeral{0}, \OLId{latin-ordinary}{z}, \OLId{latin-ordinary}{n}} \land {}}\\
      \bforall{\OLId{latin-ordinary}{i}<(\len{\OLId{latin-ordinary}{s}} \tsub \OLMathNumeral{1})}{(\OLId{latin-ordinary}{n}=\OLMathNumeral{0} \lor
      \fn{DischargeLabel}((\OLId{latin-ordinary}{s})_\OLId{latin-ordinary}{i}) \neq \OLId{latin-ordinary}{n}))}).
  \end{multline*}
  % تصحيح مصدر (R287-03): قُيّد محمول الفرض المفتوح أيضًا برمز شجرة قانوني.
  % تصحيح مصدر (C287-01، C287-02): فُصل الوسم الصفري عن أوسمة التفريغ، واستُبدل اختبار العضوية بمسار مرتب ذي حد مبرهن.
\end{proof}

\begin{prop}
  \ollabel{prop:prf-prim-rec}
  لتكن $\OLId{greek-variable}{Gamma}$ مجموعة !!{sentence}s عودية بدائية. فتكون
  العلاقة $\Prf[\OLId{greek-variable}{Gamma}](\OLId{latin-ordinary}{x},\OLId{latin-ordinary}{y})$، ومعناها «$\OLId{latin-ordinary}{x}$ رمز
  !!a{derivation}~$\delta$ لـ$!A$ من فروض !!{undischarged} تنتمي إلى
  $\OLId{greek-variable}{Gamma}$، وأن $\OLId{latin-ordinary}{y}$ عدد غودل لـ$!A$»، عودية بدائية.
\end{prop}

\begin{proof}
  ليكن المحمول العودي البدائي~$\OLId{latin-looped}{R}_\OLId{greek-variable}{Gamma}(\OLId{latin-ordinary}{y})$
  تعبيرًا عن «$\OLId{latin-ordinary}{y}\in\OLId{greek-variable}{Gamma}$».
  فنثبت عودية $\Prf[\OLId{greek-variable}{Gamma}](\OLId{latin-ordinary}{x},\OLId{latin-ordinary}{y})$
  البدائية: أي اختبار كون $\OLId{latin-ordinary}{y}$ عدد غودل
  لجملة~$!A$، وكون $\OLId{latin-ordinary}{x}$ رمز
  !!{derivation} بالاستنباط الطبيعي،
  ال!!{formula} الختامية فيه~$!A$،
  وفروضه التي هي !!{undischarged}
  منتمية كلها إلى~$\OLId{greek-variable}{Gamma}$.

  تعطي \olref{prop:deriv} أن $\fn{Deriv}(\OLId{latin-ordinary}{x})$
  عودية بدائية، ومفادها أن $\OLId{latin-ordinary}{x}$ عدد غودل
  ل!!{derivation} صحيح~$\delta$ في الاستنباط الطبيعي.
  ولذلك نعرّف $\Prf[\OLId{greek-variable}{Gamma}](\OLId{latin-ordinary}{x},\OLId{latin-ordinary}{y})$ هكذا:
  \begin{align*}
    \Prf[\OLId{greek-variable}{Gamma}](\OLId{latin-ordinary}{x}, \OLId{latin-ordinary}{y}) \defiff {}
    & \fn{Deriv}(\OLId{latin-ordinary}{x}) \land \fn{EndFmla}(\OLId{latin-ordinary}{x}) = \OLId{latin-ordinary}{y} \land {} \\
    & \bforall{\OLId{latin-ordinary}{z} < \OLId{latin-ordinary}{x}}{(\fn{OpenAssum}(\OLId{latin-ordinary}{z}, \OLId{latin-ordinary}{x}) \lif \OLId{latin-looped}{R}_\OLId{greek-variable}{Gamma}(\OLId{latin-ordinary}{z}))}.
  \end{align*}
\end{proof}

\end{document}
